\documentclass[11pt]{article}


\usepackage[margin=1in]{geometry}
\usepackage{amsmath, amssymb, amsthm}
\usepackage{enumitem}
\usepackage{booktabs}
\usepackage{graphicx}
\usepackage{xcolor}
\usepackage{listings}
\usepackage{hyperref}
\usepackage{fancyhdr}

% ---------- Header / footer ----------
\pagestyle{fancy}
\fancyhf{}
\lhead{CS4787/5777 --- Problem Set 4}
\rhead{\thepage}
\setlength{\headheight}{14pt}

% ---------- Problem / part helpers ----------
\newcommand{\problem}[2]{\section*{Problem #1: #2}}
\newcommand{\probpart}[1]{\subsection*{#1}}
\newcommand{\answer}{\par\noindent\textbf{Solution.}\ }

% ---------- Code listings ----------
\lstset{
  basicstyle=\ttfamily\small,
  breaklines=true,
  frame=single,
  columns=fullflexible
}

\title{Problem Set 4: Parallelism and Memory \\
  \large CS4787/5777 --- Principles of Large-Scale ML Systems}
\author{
  Bodong Liu - bl576 \\
  Boaz Ng - bn229 \\ % add group members (up to four) below
  Dobbs Conrad - djc478 \\
  Matthew Jia - mcj57
  }
\date{Due November 2, 2026}

\begin{document}
\maketitle

% ============================================================
\problem{1}{Parallelism and Parallel Hardware}

\probpart{1(a)}
\answer

\probpart{1(b)}
\begin{enumerate}[label=(\roman*)]
  \item \textbf{Antonio.}
  \item \textbf{Banquo.}
  \item \textbf{Cordelia.}
  \item \textbf{Desdemona.}
  \item \textbf{Edmund.}
  \item \textbf{Feste.}
  \item \textbf{Goneril.}
\end{enumerate}

\probpart{1(c)}
\answer

\probpart{1(d)}
\begin{enumerate}[label=(\roman*)]
  \item \textbf{Iago.}
  \item \textbf{Juliet.}
  \item \textbf{Kent.}
  \item \textbf{Lear.}
  \item \textbf{Miranda.}
\end{enumerate}

\newpage
% ============================================================
\problem{2}{Memory}

\probpart{2(a)}
\begin{enumerate}[label=\roman*.]
  \item Single 1024-length input token sequence:
  \item Embeddings for this token sequence:
  \item Entire dataset of input tokens:
  \item Weights of a single attention block:
  \item Weights of a single MLP block:
  \item Weights of the whole model:
  \item KV cache for a single 1024-length input sequence:
\end{enumerate}

% Optional summary table
\begin{center}
\begin{tabular}{@{}lll@{}}
\toprule
Item & Bytes & Smallest cache \\
\midrule
i.   &  &  \\
ii.  &  &  \\
iii. &  &  \\
iv.  &  &  \\
v.   &  &  \\
vi.  &  &  \\
vii. &  &  \\
\bottomrule
\end{tabular}
\end{center}

\newpage
% ============================================================
\problem{3}{Low-Precision Arithmetic}


\begin{table}[h]
\centering
\begin{tabular}{lcccc}
\toprule
\textbf{Format} & \textbf{Sign} & \textbf{Exponent} & \textbf{Mantissa} & \textbf{Total Bits} \\
\midrule
bfloat16      & 1 & 8  & 7  & 16 \\
IEEE 16       & 1 & 5  & 10 & 16 \\
MichaelFloat  & 1 & 9  & 6  & 16 \\
IEEE 32       & 1 & 8  & 23 & 32 \\
IEEE 64       & 1 & 11 & 52 & 64 \\
\bottomrule
\end{tabular}
\caption{Bit allocation for each floating-point format.}
\end{table}

\probpart{3(a) Machine epsilon}


Machine epsilon depends only on the number of mantissa bits $m$. Exponent bits only affect range. Therefore, more mantissa bits means smaller $\varepsilon_{\text{machine}}$, giving the ordering:

\[
\varepsilon_{\text{MichaelFloat}} > \varepsilon_{\text{bfloat16}} > \varepsilon_{\text{IEEE 16}} > \varepsilon_{\text{IEEE 32}} > \varepsilon_{\text{IEEE 64}}
\]

\probpart{3(b) Numerical range}

Numerical range is determined by the number of exponent bits $k$, since the
exponent sets how large or small a number can be. Mantissa bits only slightly
affect the maximum value. Therefore, more exponent bits means a larger range,
giving the ordering:
\[
\text{IEEE 64} > \text{MichaelFloat} > \text{IEEE 32} > \text{bfloat16} > \text{IEEE 16}
\]

\probpart{3(c) Quantization error}



While MichaelFloat has the fewest mantissa bits and therefore the largest machine epsilon, in terms of ML applications, data is noisy and small rounding errors typically don't matter. What ML can't recover from is a \texttt{NaN} or \texttt{inf}. Overflow produces \texttt{inf}, which quickly leads to \texttt{NaN}s, and these propagate through the model since any computation involving a \texttt{NaN} results in a \texttt{NaN}. Because of this, \fbox{we expect IEEE 16 to cause the worst quantization error} in ML applications: it has only 5 exponent bits, giving a maximum value of about $6.5 \times 10^4$, which common ML computations can exceed. MichaelFloat, by contrast, has 9 exponent bits and more range than even Float32, so overflow is not a concern.

\end{document}
